
\[
  \underline{M}\xrightarrow{\sim} f_*(\widetilde{\underline{M}})=\bar f_*(\widetilde{\underline{M}})
\]
donne par adjonction un homomorphisme~:
\[
  \bar f^{\,*}(\underline{M})\longrightarrow \underline{M}
\]
et on vérifie immédiatement sur chaque fibre que c'est un isomorphisme.

2.3.2. Soient maintenant \(\underline{B}\) une \(\underline{O}_S\)-Algèbre graduée quasi-cohérente, \(\underline{M}\) un \(\underline{B}\)-Module gradué, et \(X=\Proj(\underline{B})\). Pour tout ouvert affine \(U\) de \(S\), et toute section \(s\) de \(\underline{B}\) sur \(U\), homogène de degré \(d>0\), on note \(\underline{M}_{(s),U}\) la composante homogène de degré \(0\) du Module de fraction \(\underline{M}_s\); on a donc :
\[
  \underline{M}_{(s),U}=
  \bigl((\underline{M}|U)\otimes_{(\underline{B}|U)}(\underline{B}|U)_s\bigr)_0 .
\]
Si \(V\) est un ouvert affine de \(U\), alors \(\underline{M}_{(s),V}\) est la restriction à \(V\) de \(\underline{M}_{(s),U}\).

On considère la base d'ouverts de \(X\) formée des ouverts affines :
\[
  D_+(U,s)\simeq \Spec\bigl((\underline{B}|U)_{(s)}\bigr) .
\]
On note \(f\) le morphisme structural \(X\to S\), et \(f_{U,s}\) le morphisme affine :
\[
  f_{U,s}:D_+(U,s)\longrightarrow U .
\]
On pose alors :
\[
  \Projz(\underline{M})_{U,s}
  =
  \bar f_{U,s}^{\,*}\bigl(\underline{M}_{(s),U}\bigr),
\]
avec les notations de 2.3.1. Nous allons montrer que les \(\Projz(\underline{M})_{U,s}\) se recollent quand \(U\) et \(s\) varient, définissant ainsi un \(\underline{O}_X\)-Module noté \(\Projz(\underline{M})\). Le recollement pour \(U\) variable et \(s\) fixe étant évident, on est ramené (cf. EGA II 2.5.2) à montrer que si \(t\) est une section de \(\underline{B}\) sur \(U\), homogène de degré \(e>o\), alors la restriction de \(\Projz(\underline{M})_{U,s}\) à \(D_+(U,st)\) est canoniquement isomorphe à \(\Projz(\underline{M})_{U,st}\). On considère :
